\section{Spectral triples, zeta cycles, and prolate operators}
\label{sec:ACspectral}

\subsection{Sources and scope}
In this section \(A\) is the input theorem of the question.  Its zeta
component is used; the Dirichlet component is not needed here.  The source
abbreviations are as follows.  \(\mathrm{S21}\) is Connes--Consani,
\emph{Spectral triples and zeta-cycles}, arXiv:2106.01715, attached
\texttt{Spectraltriples.tex}.  \(\mathrm{S22}\) is their
\emph{Hochschild homology, trace map and zeta-cycles}, arXiv:2207.10419,
attached \texttt{Hochschild-zeta.tex}.  \(\mathrm{S23}\) is
Connes--Consani--Moscovici, \emph{Zeta zeros and prolate wave operators},
arXiv:2310.18423, attached \texttt{mainc2m24fine.tex}.
\(\mathrm{S25}\) is their \emph{Zeta spectral triples},
arXiv:2511.22755, attached \texttt{mc2arXiv.tex}.  Source labels below are
literal LaTeX labels, so they remain exact locators if the numbering changes
between versions.  The source numbering used for S21, S22 and S25 is that
of the supplied text.

The following distinctions are essential.  A spectrum contained in the
real line does not by itself give a positive inner product making a
specified operator self-adjoint.  A positive quadratic form does not by
itself identify its eigenvectors.  A counting-function asymptotic does
not identify individual spectral points.  The results below give exact
positive conclusions and identify the additional hypotheses at each of
these three transitions.

\subsection{What becomes complete in the zeta-cycle and sheaf realizations}

\begin{proposition}[AC-S01: all zeros are covered]
\label{AC-S01}
Let
\[
 C_L=\mathbb R_+^*/\exp(L\mathbb Z),\qquad
 \mathcal H_L=\overline{\Sigma_{e^L}\mathcal E(\mathcal S_0^{\rm ev})}^{\,\perp}
 \subset L^2(C_L),\quad L>0.
\]
Under \(A\), the spectrum of the self-adjoint generator of scaling on
\(\mathcal H_L\), with the sign convention in which a character is
\(u^{it}\), is exactly
\[
 \{\gamma\in\mathbb R:\zeta(1/2+i\gamma)=0,
                       \ \gamma L\in2\pi\mathbb Z\}.
\]
Each character in a single circle has multiplicity one.  Every nontrivial
zeta zero occurs for some \(L\), and the set of such lengths is countable
and dense in \(\mathbb R_{>0}\).  The sheaf cohomology realization of
S22, Theorems 1.2 and 5.4, now covers all nontrivial zeros, with their
local jet multiplicities.
The Laplacian realization in S22, Proposition 4.2 and Corollary 4.3
(labels \texttt{laplacian}, \texttt{rhequiv}), has spectrum
\[
 \{-1/4-\gamma^2:\zeta(1/2+i\gamma)=0\},
\]
with the multiplicity conventions of that proposition.
\end{proposition}
\begin{proof}
S21, Theorem 6.4 and Corollary 6.5 (labels \texttt{spectralreal},
\texttt{repetitions}), identify the orthogonal complement by precisely
the critical zeros whose characters are periodic on the indicated circle.
The Fourier modes of \(L^2(C_L)\) are one-dimensional, which proves the
single-circle multiplicity statement.  Theorem \(A\) identifies critical
zeros with all nontrivial zeros.  For any positive zero ordinate
\(\gamma\), choose \(L=2\pi n/\gamma\), \(n\geq1\).
Conversely every nonzero \(\mathcal H_L\) has this form for some
\(\gamma,n\).  This is a countable set.  For any \(L_0>0\), take
zero ordinates \(\gamma_j\to\infty\), and choose positive integers
\(n_j\) nearest to \(L_0\gamma_j/(2\pi)\).  Then
\(2\pi n_j/\gamma_j\to L_0\), proving density.
S22, Theorem 5.4, identifies its cohomology with the Schwartz quotient
by the closed ideal generated by \(t\mapsto\zeta(1/2+it)\); its local
description records the zero jets.  Finally substitute
\(\rho=1/2+i\gamma\) into the eigenvalue
\((\rho-1/2)^2-1/4\) of S22, Proposition 4.2.
\end{proof}

\begin{proposition}[AC-S02: the multiplicity obstruction to unitarization]
\label{AC-S02}
At a zero of order \(m\) and real ordinate \(\gamma\), the local jet
representation in S22 is the module
\(J_\gamma=\mathbb C[t]/(t-\gamma)^m\).  On this module, multiplication
by \(t\) is self-adjoint for some positive definite Hermitian form if
and only if \(m=1\).  Equivalently, the scaling representation on this
finite-dimensional block is unitarizable if and only if \(m=1\).
For \(\gamma\ne0\), the same obstruction applies to the Laplacian
\(-1/4-t^2\).  Thus AC-S01 proves a statement about its spectrum;
unitarizing the given jet realization additionally requires simplicity
of the zeros.  Theorem \(A\) as stated does not supply that hypothesis.
\end{proposition}
\begin{proof}
In the basis \(1,t-\gamma,\ldots,(t-\gamma)^{m-1}\), multiplication
by \(t\) is \(\gamma I+J\), with \(J\) a nilpotent Jordan block of
length \(m\).  If \(m>1\), it is not diagonalizable.  Every
self-adjoint operator in a positive finite-dimensional inner product is
diagonalizable, so no such inner product exists.  A unitary one-parameter
group has a self-adjoint infinitesimal generator, giving the equivalent
formulation.  If \(m=1\), the usual Hermitian form works.  The
nilpotent part of \(-1/4-(\gamma I+J)^2\) is
\(-2\gamma J-J^2=-J(2\gamma I+J)\).  If \(\gamma\ne0\), the
second factor is invertible, so for \(m>1\) this has nonzero nilpotent
part and again is not diagonalizable.  There is no zeta zero at ordinate
zero: \(\zeta(1/2)<0\), for example by its alternating-series
continuation.  The argument is an exact obstruction for multiple zeros,
not a claim that a multiple zero exists.
\end{proof}

\subsection{The closed Weil operators on finite windows}

Write \(L=\log\lambda\), and identify
\(L^2([\lambda^{-1},\lambda],d^*u)\) with
\(H_L=L^2([-L,L],dx)\).  All functions on the interval are extended by
zero.  Use
\(\widehat f(t)=\int_{\mathbb R}f(x)e^{-itx}\,dx\).
Let \(q_L\) be the closed lower semibounded Weil form of S25,
\S3.1, and let \(A_L\) be its associated operator.
S25, Proposition 3.4 (\texttt{Hilbert1}), gives the Fourier-polynomial
form core; Theorem 3.6 (\texttt{thmsmallest}) gives compact resolvent.
Write its eigenvalues with multiplicity as
\(a_1(L)\leq a_2(L)\leq\cdots\).

\begin{theorem}[AC-S03: positive gap and exact Galerkin limits]
\label{AC-S03}
Under \(A\), for every finite \(L>0\),
\[
 q_L(f)=\sum_{\gamma}m_\gamma|\widehat f(\gamma)|^2
 \quad(f\in\operatorname{Dom}q_L),\qquad
 0<a_1(L)\leq a_2(L)\leq\cdots\longrightarrow\infty.
\]
Here the sum is over distinct ordinates of all nontrivial zeta zeros.
For the nested Fourier-polynomial spaces \(E_N\) of S25, let
\(a_{r,N}(L)\) denote the \(r\)-th eigenvalue of the restricted form.
For every fixed \(r\), once \(\dim E_N\geq r\),
\[
 a_{r,N}(L)\downarrow a_r(L)>0.
\]
The result is strict positivity of the Weil operator and its finite
matrices; it makes no assertion that \(a_1(L)\) is simple or even.
Strict positivity under RH is already stated in the introduction of
S21.  The proof here supplies the full closed-domain and spectral-gap
details and the higher Galerkin limits.
\end{theorem}
\begin{proof}
The explicit formula and \(A\) give the displayed sum on the
piecewise smooth core.  For \(f\in\operatorname{Dom}q_L\), choose
core vectors \(f_n\to f\) in form norm.  The vectors
\((\sqrt{m_\gamma}\widehat f_n(\gamma))_\gamma\) are Cauchy in
\(\ell^2\), since their squared distance is \(q_L(f_n-f_k)\).
For each fixed \(\gamma\), evaluation of the Fourier transform is
continuous on \(H_L\), by
\(|\widehat f(\gamma)|\leq\sqrt{2L}\|f\|_2\).
Consequently the limit has exactly the displayed coordinates, and the
identity extends to the closed form domain.

If \(q_L(f)=0\), its entire Fourier transform vanishes at every
distinct ordinate.  A nonzero Fourier transform of an \(L^2\) function
supported in \([-L,L]\) has only \(O_f(T)\) zeros in \(|z|\leq T\),
counting multiplicity.  Indeed
\(|\widehat f(z)|\leq\sqrt{2L}\|f\|_2e^{L|\operatorname{Im}z|}\),
and Jensen's formula about any point at which the transform is nonzero,
applied on discs of radii \(T\) and \(2T\), gives that bound.
The distinct zeta ordinates have superlinear counting growth.  One may
use Conrey's unconditional positive-proportion simple-zero theorem
(\emph{Bull. Amer. Math. Soc.} 20 (1989), Theorem 1, p.~79), together
with the Riemann--von Mangoldt formula.  Alternatively, directly under
\(A\), Carneiro--Chandee--Milinovich,
\emph{Bounding \(S(t)\) and \(S_1(t)\) on the Riemann hypothesis},
arXiv:1309.1526, Theorem 2 and equation (1), give
\(m_\gamma=O(\log T/\log\log T)\) on \(T/2\leq\gamma\leq T\)
and \(N(T)-N(T/2)\gg T\log T\).  The number of distinct ordinates in
that interval is therefore \(\gg T\log\log T\).
Both routes contradict the Jensen bound unless \(f=0\).
The compact-resolvent theorem cited above means that the bottom of the
spectrum is an attained eigenvalue.  Its nonnegativity and absence of a
kernel therefore imply \(a_1(L)>0\).

The minimum--maximum principle gives
\(a_{r,N+1}(L)\leq a_{r,N}(L)\) and
\(a_{r,N}(L)\geq a_r(L)\).  Approximate an orthonormal basis of the
span of the first \(r\) eigenvectors in form norm by vectors in a
single \(E_N\).  On this fixed finite-dimensional span the Gram and
form matrices converge, so the largest Rayleigh quotient converges to
\(a_r(L)\).  Applying the same principle gives the reverse limit
inequality.  This proves the claimed convergence.
\end{proof}

\begin{lemma}[AC-S04: a normalized global null function]
\label{AC-S04}
Define
\[
 h(u)=2\pi u^2(2\pi u^2-3)e^{-\pi u^2},\qquad
 K(x)=e^{x/2}\sum_{n\geq1}h(ne^x).
\]
Then \(K\) is a nonzero even smooth function, every derivative decays
faster than every exponential at both ends, and
\[
 \widehat K(t)=\Xi(t):=\xi(1/2+it),\qquad
 \xi(s)=\tfrac12s(s-1)\pi^{-s/2}\Gamma(s/2)\zeta(s).
\]
For every integer \(j\geq0\) and \(u=e^x\geq2\),
\[
 |K^{(j)}(x)|\leq C_j u^{9/2+2j}e^{-\pi u^2}.
\]
With the unnormalized Fourier transform explicitly used in S25,
equation \texttt{mhat}, the function \(h\) printed in equation
\texttt{ku} of that source is one quarter of the one above.  Its
transform is \(\Xi/4\), so the equality as printed in Lemma 7.1
(\texttt{hermfact}) requires this scalar correction.
\end{lemma}
\begin{proof}
Gaussian Fourier differentiation gives
\(\mathcal F(h)=h\) for the Fourier kernel \(e^{2\pi iut}\).
Explicitly, the transforms of \(u^2e^{-\pi u^2}\) and
\(u^4e^{-\pi u^2}\) are respectively
\((1/(2\pi)-t^2)e^{-\pi t^2}\) and
\((t^4-3t^2/\pi+3/(4\pi^2))e^{-\pi t^2}\); substituting the
coefficients of \(h\) proves the assertion.  Also \(h(0)=0\) and
\(\int h=\mathcal Fh(0)=0\).  Poisson summation, pairing positive
and negative integers, therefore gives \(K(x)=K(-x)\).
Termwise differentiation on \(u\geq2\) is justified by Gaussian
domination.  The \(j\)-th derivative is a sum of
\(u^{1/2}P_j(nu)e^{-\pi n^2u^2}\), where \(P_j\) has degree at
most \(4+2j\).  Factoring out \(e^{-\pi u^2}\) bounds the remaining
sum uniformly for \(u\geq2\), and proves the stated estimate and
the decay at both ends.

For \(\operatorname{Re}s>1\), Fubini and the substitution \(v=nu\)
give
\[
 \int_0^\infty u^{1/2}\sum_{n\geq1}h(nu)\,
                  u^{s-1/2}\frac{du}{u}
   =\zeta(s)\int_0^\infty h(v)v^{s-1}\,dv.
\]
The elementary Gaussian integral gives
\[
 \int_0^\infty h(v)v^{s-1}\,dv
   =\tfrac12s(s-1)\pi^{-s/2}\Gamma(s/2).
\]
The rapid decay already proved makes the integral on the left entire
in \(s\).  Analytic continuation, followed by \(s=1/2-it\) and
\(\xi(s)=\xi(1-s)\), proves the Fourier identity.  Since \(\xi\)
is not identically zero, \(K\ne0\).  Repeating the same integral
with the coefficient \(\pi/2\) of S25 instead of \(2\pi\) divides
the result by four, proving the scalar correction.
\end{proof}

\begin{theorem}[AC-S05: every fixed low Weil eigenvalue tends to zero]
\label{AC-S05}
Under \(A\), for each fixed positive integer \(r\) there are
\(C_r<\infty\) and \(\lambda_r\) such that
\[
 0<a_r(\log\lambda)
 \leq C_r\lambda^{4r+5}e^{-2\pi\lambda^2}
 \quad(\lambda\geq\lambda_r).
\]
In particular \(\mu_\lambda=a_1(\log\lambda)\) is positive for every
finite window, decreases to zero, and satisfies
\(\mu_\lambda=O(\lambda^9e^{-2\pi\lambda^2})\).
Together with S25, Corollary 3.8 (\texttt{strange}), this gives
\[
 \mathrm{RH}\quad\Longleftrightarrow\quad
 \lim_{\lambda\to\infty}\mu_\lambda=0.
\]
It does not give the numerical asymptotic
\(\mu_\lambda\asymp e^{-4\pi\lambda^2}\) discussed in S25.
\end{theorem}
\begin{proof}
Put \(L=\log\lambda\), \(K_j=K^{(j)}\), and
\(f_{j,L}=1_{[-L,L]}K_j\), \(0\leq j<r\).
These are piecewise smooth compactly supported functions and belong
to the form domain; this is also explicitly observed immediately after
S25, Corollary 3.7.  They need not lie in the operator domain.
Here is a verification in the full form norm.  For fixed \(L\), multiply
\(K_j\) by a smooth cutoff equal to one on
\([-L+\delta,L-\delta]\) and supported in \((-L,L)\), with derivative
bounded by \(C/\delta\).  The difference \(w_\delta\) from
\(f_{j,L}\) is supported in two intervals of total length \(2\delta\),
is uniformly bounded, and has uniformly bounded variation, including
its endpoint jumps.  Hence
\[
 |\widehat w_\delta(t)|\leq C_{j,L}
                   \min\{\delta,|t|^{-1}\}.
\]
The archimedean symbol is bounded in absolute value by
\(C_L+\log(2+|t|)\); the finite prime and pole terms are bounded
operators on \(H_L\), as in S25, the proof of Theorem 3.6.
Splitting the Fourier integral at \(|t|=\delta^{-1}\) gives
\[
 \|w_\delta\|_2^2+
 \int_{\mathbb R}\log(2+|t|)|\widehat w_\delta(t)|^2\,dt
       =O_{j,L}(\delta\log(1/\delta)).
\]
Thus the smooth cutoffs converge to \(f_{j,L}\) in the closed form
norm.  In particular, membership in the form domain and the spectral
identity used below are justified by form convergence, not just by
\(L^2\) convergence.
The Fourier transforms of the uncut functions are
\((it)^j\Xi(t)\), and vanish at every zero ordinate.
For \(R_{j,L}=K_j-f_{j,L}\), the derivative in the sense of measures
has two boundary atoms as well as the derivatives on the two tails.
The estimates of AC-S04 and
\[
 \int_\lambda^\infty u^a e^{-\pi u^2}\frac{du}{u}
       \leq C_a\lambda^{a-2}e^{-\pi\lambda^2}
 \quad(\lambda\geq\lambda_a)
\]
show that its total variation satisfies
\[
 \operatorname{Var}(R_{j,L})
 \leq C_j\lambda^{9/2+2j}e^{-\pi\lambda^2}.
\]
Integration by parts against that finite measure gives, for \(t\ne0\),
\[
 |\widehat R_{j,L}(t)|
       \leq |t|^{-1}\operatorname{Var}(R_{j,L}).
\]
The Riemann--von Mangoldt estimate implies
\(B:=\sum_\gamma m_\gamma/\gamma^2<\infty\); there is no zero
ordinate at zero.  For \(a=(a_0,\ldots,a_{r-1})\), the spectral
formula of AC-S03 and the vanishing of \(\widehat K_j(\gamma)\)
therefore give
\[
 q_L\left(\sum_{j<r}a_jf_{j,L}\right)
 \leq C_r\|a\|^2\lambda^{4r+5}e^{-2\pi\lambda^2}.
\]
The functions \(K_0,\ldots,K_{r-1}\) are linearly independent:
after Fourier transform a linear relation is
\(P(it)\Xi(t)=0\), which forces the polynomial \(P\) to vanish.
Their Gram matrix is thus positive definite.  The Gram matrices of
the cutoffs converge to it in every entry by dominated convergence.
Consequently, for all sufficiently large \(L\),
\[
 \left\|\sum_{j<r}a_jf_{j,L}\right\|_2^2
           \geq c_r\|a\|^2\quad(c_r>0).
\]
This is an \(r\)-dimensional trial space.  The minimum--maximum
principle proves the upper bound; AC-S03 proves strict positivity.
Monotonicity for \(r=1\) follows by extending test functions by zero
as the interval grows, or from S25, equation \texttt{monotone}.
The converse in the final equivalence is exactly S25, Corollary 3.8.

An explicit constant is available for the lowest eigenvalue.  Put
\(c_0=\int_{-\log2}^{\log2}|K(x)|^2\,dx>0\) and
\(B=\sum_\gamma m_\gamma/\gamma^2\).  For \(u\geq2\), every
summand of \(K(\log u)\) is positive and decreasing: its logarithmic
derivative has the sign of
\(-\pi v^2(8\pi^2v^4-30\pi v^2+15)\), \(v=nu\).
Consequently \(\operatorname{Var}(R_{0,L})=4K(L)\).
Also
\[
 K(\log\lambda)\leq
 4\pi^2\lambda^{9/2}e^{-\pi\lambda^2}
       \sum_{n\geq1}n^4e^{-4\pi(n^2-1)}
 \leq8\pi^2\lambda^{9/2}e^{-\pi\lambda^2}
 \quad(\lambda\geq2).
\]
For the last bound, use \(n^2-1\geq n-1\) and
\(\sum_{n\geq1}n^4q^{n-1}=(1+11q+11q^2+q^3)/(1-q)^5<2\)
at \(q=e^{-4\pi}\).  Therefore the concrete inequality
\[
 \mu_\lambda\leq
 \frac{1024\pi^4B}{c_0}\lambda^9e^{-2\pi\lambda^2}
 \quad(\lambda\geq2)
\]
holds.  Both \(B\) and \(c_0\) are fixed positive constants defined
by absolutely convergent expressions.
\end{proof}

\subsection{The convergence that can and cannot presently be asserted}

\begin{proposition}[AC-S06: fixed-window eigenfunction convergence]
\label{AC-S06}
Fix \(L>0\).  Any sequence of \(L^2\)-normalized lowest Galerkin
eigenvectors has a subsequence converging in \(L^2\) and in form norm
to a lowest eigenvector of \(A_L\).  Their Fourier transforms converge
uniformly on compact subsets of \(\mathbb C\) along that subsequence.
If the lowest eigenvalue of \(A_L\) is simple, phases can be chosen so
that the whole sequence converges.  If that eigenvector is also even,
then for all sufficiently large \(N\) the Galerkin matrix is even-simple,
and the normalized determinant of S25, Theorem 5.10
(\texttt{finmain}), converges locally uniformly to a nonzero scalar
multiple of the Fourier transform of that eigenvector, after removing
the fixed phase factor \(-i\lambda^{-iz}\).
This assertion permits an \(L^2\) normalization.  It does not assert
convergence of the boundary normalization \(\xi_N(\lambda)=1\).
\end{proposition}
\begin{proof}
The eigenvalues converge by AC-S03.  Boundedness of the form norm and
compactness of its embedding in \(H_L\) give a strongly \(L^2\)
convergent subsequence and a weakly form-convergent subsequence with
the same limit.  For \(g\) in the union of the \(E_N\), the Galerkin
eigenvalue identity eventually holds.  Passing to the limit gives
\(q_L(f,g)=a_1(L)\langle f,g\rangle\).  The form-core property
extends this equality to the whole form domain.  Hence \(f\) is a
normalized lowest eigenvector.  Its form norm equals the limiting form
norm of the approximants, so weak convergence in that Hilbert norm is
strong.  Simplicity identifies every subsequential limit up to phase.

For a compact set \(|\operatorname{Im}z|\leq R\),
\[
 |\widehat f_N(z)-\widehat f(z)|
 \leq\left(\int_{-L}^Le^{2R|x|}\,dx\right)^{1/2}\|f_N-f\|_2,
\]
proving local uniform convergence.  If the lowest eigenvector is simple
and even, the spectral gap to the second eigenvalue and convergence
of both Galerkin eigenvalues make the first Galerkin eigenvalue simple
for large \(N\).  Inversion preserves \(E_N\), and its eigenvectors
have parity.  Convergence to the even lowest vector forces even parity.
The boundary functional on that finite-dimensional eigenvector is
nonzero by the argument between Definition 5.3 and Lemma 5.4 of S25.
The determinant formula of Theorem 5.10 is thus applicable.  Replacing
its boundary normalization by the \(L^2\) normalization multiplies
its determinant formula by a nonzero constant; the preceding transform
convergence proves the claim.  Point evaluation at a boundary is not
continuous in \(L^2\), so no boundary-normalization limit was used.
\end{proof}

\begin{proposition}[AC-S07: a precise sufficient gate for the large-window limit]
\label{AC-S07}
Let \(u_L\) be a normalized even simple lowest eigenvector, and let
\(v_L\) be a normalized trial vector in \(\operatorname{Dom}q_L\).
Write \(g_L=a_2(L)-a_1(L)>0\).  After choosing a phase,
\[
 \|u_L-v_L\|_2^2
 \leq 2\frac{q_L(v_L)-a_1(L)}{g_L}.
\]
If trial vectors \(c_Lv_L\), extended by zero, converge to \(K\)
in \(L^1(e^{b|x|}dx)\), and
\[
 |c_L|^2\left(\int_{-L}^Le^{2b|x|}\,dx\right)
       \frac{q_L(v_L)-a_1(L)}{g_L}\longrightarrow0,
\]
then \(c_L\widehat u_L\to\Xi\) uniformly on
\(|\operatorname{Im}z|\leq b\).  The same statement with a compact
strip region follows from its corresponding weaker norm bound.
\end{proposition}
\begin{proof}
Expand \(v_L=\alpha u_L+w\), \(w\perp u_L\), and choose the
phase so that \(\alpha\geq0\).  The spectral theorem for the closed
form gives
\(q_L(v_L)-a_1(L)\geq g_L\|w\|^2\), while
\(\|v_L-u_L\|^2=2(1-\alpha)\leq2(1-\alpha^2)=2\|w\|^2\).
Weighted Cauchy--Schwarz on \([-L,L]\) bounds the weighted \(L^1\)
distance between \(c_Lu_L\) and \(c_Lv_L\) by the square root of
the second displayed expression, times \(\sqrt2\).
The triangle inequality and
\(|\widehat f(z)|\leq\int|f(x)|e^{b|x|}dx\) complete the proof.
\end{proof}

\begin{proposition}[AC-S13: the prolate trial function has a proved strip limit]
\label{AC-S13}
Let \(h\) be the correctly normalized function in AC-S04.  Suppose
\(h_\lambda\) is supported in \([-\lambda,\lambda]\) and
\[
 \sup_{|u|\leq\lambda}|h_\lambda(u)-h(u)|\leq C\lambda^{-2}.
\]
This is the estimate in S25, Lemma 7.2 (\texttt{meixnerlem}), after
the same scalar normalization as AC-S04.  Define
\[
 k_\lambda(x)=1_{[-\log\lambda,\log\lambda]}(x)e^{x/2}
                        \sum_{n\geq1}h_\lambda(ne^x).
\]
For every \(0\leq b<1/2\),
\[
 \|k_\lambda-K\|_{L^1(e^{b|x|}dx)}
      =O_b(\lambda^{b-1/2}),\qquad
 \sup_{|\operatorname{Im}z|\leq b}
          |\widehat k_\lambda(z)-\Xi(z)|
      =O_b(\lambda^{b-1/2}).
\]
This is an unconditional trial-function limit.  It does not assert
that \(k_\lambda\) is a ground eigenvector.
\end{proposition}
\begin{proof}
Write \(u=e^x\), \(L=\log\lambda\), and
\(w_b(u)=\max(u^b,u^{-b})\).  On \(\lambda^{-1}\leq u\leq\lambda\),
the at most \(\lambda/u\) retained terms contribute an error at most
\(C\lambda^{-1}u^{-1/2}\).  Its weighted integral is bounded by
\[
 C\lambda^{-1}\int_{\lambda^{-1}}^\lambda
          u^{-1/2}w_b(u)\frac{du}{u}
       =O_b(\lambda^{b-1/2}).
\]
There is also a tail from the terms of \(h\) with \(nu>\lambda\),
which is not included in that first estimate.  For \(\lambda\geq2\),
\(h\) is positive and decreasing on \([\lambda,\infty)\).
An integral comparison bounds this tail by
\[
 u^{1/2}\sum_{nu>\lambda}|h(nu)|
 \leq u^{1/2}h(\lambda)+u^{-1/2}\int_\lambda^\infty h(v)\,dv.
\]
Here \(h(\lambda)=O(\lambda^4e^{-\pi\lambda^2})\) and
\(\int_\lambda^\infty h(v)\,dv=O(\lambda^3e^{-\pi\lambda^2})\).
Its weighted integral on the window is therefore
\(O_b(\lambda^{b+9/2}e^{-\pi\lambda^2})\).
Finally the weighted integral of \(K\) outside the window is
\(O_b(\lambda^{b+5/2}e^{-\pi\lambda^2})\), by its evenness and
AC-S04.  This proves the first estimate, including both omitted tails.
The Fourier estimate follows from the weighted \(L^1\) bound.
\end{proof}

The two obstacles expressly listed in S25, \S8 (\texttt{missing}),
are simplicity/evenness and comparison with the prolate candidate.
AC-S03 and AC-S05 do not settle them.  In particular, AC-S05 also
makes \(a_2(L),a_3(L),\ldots\) tend to zero; a small Rayleigh quotient
without a comparison to the spectral gap cannot select the ground
state.  AC-S07 makes the missing comparison quantitative.  Positivity
and parity invariance alone do not force an even simple minimum:
in a two-dimensional space with one even and one odd vector,
\(\operatorname{diag}(2,1)\) has an odd minimum and \(I\) has a
multiple minimum.  These are counterexamples to that abstract
inference, not counterexamples for the specific Weil matrices.

S21, \S\S4--5, exhibits numerical agreements for its different
operators \(D(\lambda,k)\); Proposition 4.2 proves a counting
asymptotic, and Theorem 6.4 proves the exact zeta-cycle realization.
Neither is a convergence theorem for \(D(\lambda,k)\) along an
unspecified joint limit.  A proof of that numerical convergence still
needs approximation estimates for the chosen prolate projection,
a specified path \(k=k(\lambda)\), and control of spectral pollution.
For S25, the fixed-window result above is a proved form of convergence;
the further \(\lambda\to\infty\) determinant convergence retains
the stated gates.  This is not a claim of logical independence from RH.
\subsection{A spectral determinant construction with a proved limit}

\begin{theorem}[AC-S08: an explicit Hilbert realization and determinant]
\label{AC-S08}
Under \(A\), form the Hilbert space with orthonormal basis
\(e_{\gamma,j}\), where \(\gamma\) ranges over distinct positive and
negative ordinates and \(1\leq j\leq m_\gamma\), and put
\[
 De_{\gamma,j}=\gamma e_{\gamma,j},\qquad
 \operatorname{Dom}D=
 \left\{(c_{\gamma,j}):\sum_{\gamma,j}\gamma^2|c_{\gamma,j}|^2<\infty\right\}.
\]
Then \(D\) is self-adjoint with compact resolvent,
\(D^{-1}\) is Hilbert--Schmidt, and
\[
 \det{}_2(I-zD^{-1})=\frac{\Xi(z)}{\Xi(0)}.
\]
For the spectral cutoff \(P_T=1_{[-T,T]}(D)\),
\[
 \left\|D^{-1}-P_TD^{-1}P_T\right\|_2^2
       =O(\log T/T),\qquad
 \det(I-zD_T^{-1})\longrightarrow\Xi(z)/\Xi(0),
\]
where \(D_T\) acts on \(P_TH\), and the determinant convergence is
locally uniform with error \(O_R(\log T/T)\) on \(|z|\leq R\).
The diagonal algebra \(c_0(\{(\gamma,j)\})\), or its unitization,
with this \(D\) gives a spectral triple; its summability threshold is
\(p>1\).  This construction uses the zero multiset and does not identify
the operators of S21 or S25 with these diagonal cutoffs.
\end{theorem}
\begin{proof}
The multiplication operator on the stated maximal domain is self-adjoint.
There are finitely many ordinates in any bounded interval, and their
absolute values tend to infinity, so its resolvent is compact.  Partial
summation of \(N(T)=O(T\log T)\) gives
\(\sum_{\gamma>T}m_\gamma/\gamma^2=O(\log T/T)\), proving both
Hilbert--Schmidt assertions.  The order-one Hadamard factorization of
the even entire function \(\Xi\), paired at \(\gamma,-\gamma\),
is
\[
 \Xi(z)/\Xi(0)=\prod_{\gamma>0}(1-z^2/\gamma^2)^{m_\gamma}.
\]
Indeed the paired product converges by the preceding summability;
Hadamard factorization leaves an exponential \(e^{az+b}\), whose
linear term vanishes by evenness and whose constant is fixed at zero.
The defining regularized determinant product
\(\prod_{\gamma,j}(1-z/\gamma)e^{z/\gamma}\) has precisely these
paired factors.  The finite ordinary determinant has the truncated
paired product.  For \(T>2R\), the logarithm of its omitted tail is
bounded by
\(C_R\sum_{\gamma>T}m_\gamma/\gamma^2\).
The truncated products are bounded uniformly on \(|z|\leq R\) by
\(\exp(R^2\sum_{\gamma>0}m_\gamma/\gamma^2)\), which gives the
stated absolute error.  Finite diagonal sequences form a dense
\(*\)-subalgebra with zero commutators with \(D\); compact resolvent
gives the remaining spectral-triple condition.  Finally partial
summation of \(N(T)\sim T\log T/(2\pi)\) proves convergence of
\(\operatorname{Tr}(1+D^2)^{-p/2}\) exactly for \(p>1\).
\end{proof}

\subsection{What the semilocal prolate construction actually supplies}

For a finite set of places \(S\ni\infty\), S23,
Propositions \texttt{groundstate} and \texttt{httransfoS}, gives the
unitary model
\[
 D_S=M_t\ \text{on }L^2(\mathbb R,dm_S),\qquad
 dm_S(t)=c_S\left|\Gamma(1/4+it/2)\right|^2
       \prod_{p\in S\setminus\{\infty\}}
          |1-p^{-1/2+it}|^{-2}\,dt,
\]
where a positive constant \(c_S\) may normalize the cyclic vector.
Let \(N_S\) be the number operator of the associated orthogonal
polynomials.  The formal candidate in Definition \texttt{prolop} is
\(T_{\lambda,S}=-D_S^2+\lambda^2N_S\); the physical archimedean
operator uses the rescaling and additive constant in S23, equation
\texttt{wlambdconj1}.

\begin{proposition}[AC-S09: spectral type and self-adjoint extensions]
\label{AC-S09}
For every finite \(S\), \(D_S\) has purely absolutely continuous
spectrum \(\mathbb R\), of multiplicity one, and no eigenvalues.
The formal \(T_{\lambda,S}\) on the finite span of orthogonal
polynomials is densely defined and symmetric.  It admits self-adjoint
extensions; extensions preserving the even/odd grading can be chosen.
Their existence and real spectra are unconditional.  The construction
does not specify a unique extension or establish a discrete zeta
spectrum for it.
\end{proposition}
\begin{proof}
The displayed density is strictly positive and finite everywhere on
\(\mathbb R\), since the gamma function has no zeros or poles on
this line and \(p^{-1/2}<1\).  Multiplication by the square root of
the density identifies this model with multiplication by \(t\) on
ordinary \(L^2(\mathbb R,dt)\).  This proves the spectral assertion.
Every polynomial has all required moments, and multiplication by
\(t^2\) and \(N_S\) preserve the polynomial span, on which they
are symmetric.  Complex conjugation \(J\) commutes with this
 symmetric operator.  The span is dense: the finite Euler factors
give norms equivalent to the archimedean norm, in which completeness
of these polynomials follows from S23, Proposition \texttt{numberN}
and its unitary transform \texttt{httransfo}.
Complex conjugation \(J\) therefore maps
\(\ker(T_{\lambda,S}^*-i)\) antiunitarily onto
\(\ker(T_{\lambda,S}^*+i)\), so the deficiency indices agree.
Explicitly, choose an orthonormal basis \((e_j)\) of the first
deficiency space and the basis \((Je_j)\) of the second, and let
\(Ue_j=Je_j\) with \(U\) extended linearly.  The extension has
domain
\[
 \operatorname{Dom}\overline T\ \dotplus\
 \{h+Uh:h\in\ker(T^*-i)\},
\]
and acts by \(\overline T f+ih-iUh\).  The standard von Neumann
extension theorem gives self-adjointness.  One may make these choices
separately on the two parity subspaces because \(T\) preserves them.
Theorem \(A\) is not used in any step.
\end{proof}

\begin{proposition}[AC-S10: the semilocal number operator is not a transported one]
\label{AC-S10}
If \(S\) contains a finite prime, the identity map on polynomials
between the archimedean and semilocal weighted spaces does not
intertwine their number operators.  Thus the bounded invertible
\(\eta_S\) of S23 cannot be used to transport the archimedean
prolate spectrum by the claimed intertwining
\(N_S\eta_S=\eta_SN_\infty\).
\end{proposition}
\begin{proof}
If it did intertwine them, every archimedean orthogonal polynomial
of degree \(n\) would also be an eigenvector of the semilocal
\(N_S\) for \(n\).  Since these eigenspaces are one-dimensional,
the two monic orthogonal-polynomial systems would coincide.
In particular the integral of every positive-degree archimedean
orthogonal polynomial against the normalized semilocal measure would
vanish.  Expanding arbitrary polynomials in that basis shows that all
normalized moments of the two measures would agree.  These moments
are determinate: the archimedean Jacobi coefficients are
\(a_n=\sqrt{(n+1)(n+1/2)}\), and
\(\sum a_n^{-1}=\infty\) gives the Carleman determinacy criterion
(S23, Proposition \texttt{fourierfmuall}, and the determinacy argument
in Theorem \texttt{dzero}, part (ii)).
But their density ratio is a constant times
\[
 R_S(t)=\prod_{p\in S\setminus\{\infty\}}
                         |1-p^{-1/2+it}|^{-2}.
\]
This function is not constant: each factor is at most its value at
\(t=0\), and choosing \(t\) for which one of the phases is not one
makes the inequality strict.  Thus the measures are different, a
contradiction.  The conclusion agrees with S23, Remark
\texttt{care}(i); it has been re-proved here to identify exactly
what prevents this transfer.
\end{proof}

\begin{proposition}[AC-S11: correction of the Jacobi metaplectic normalization]
\label{AC-S11}
On finite sequences, let \(N e_n=ne_n\), and let \(B\) be a real
Jacobi matrix with zero diagonal and positive off-diagonal entries
\(b_n\).  Put \(C=[B,N]\).
The assignments \(h\mapsto-iB\), \(k\mapsto N+c\) in S23,
\S5.2, equation \texttt{h}, are incompatible with the standard
relations \([h,e_\pm]=\pm2e_\pm\), \([e_+,e_-]=h\) and
\(k=i(e_--e_+)\): they give
\([\sigma(e_+),\sigma(e_-)]=\sigma(h)/4\).
The corrected assignments
\[
 \sigma(h)=-iB,\quad \sigma(k)=2(N+c),\quad
 \sigma(e_+)=i(N+c)+\tfrac12C,\quad
 \sigma(e_-)=-i(N+c)+\tfrac12C
\]
satisfy the Lie algebra relations if and only if
\[
 b_n^2=(n+1)(n+2c)\quad(n\geq0).
\]
For every \(c>0\) these entries are positive, so the Lie algebra
relations alone do not select \(c=1/4\).  The even metaplectic
realization selects that value by its compact-generator normalization;
it then gives the coefficients printed in S23, equation
\texttt{coefffixed}.  The uniqueness assertion of Theorem
\texttt{dzero}(i) requires this extra normalization.
\end{proposition}
\begin{proof}
Directly from the tridiagonal matrix,
\([[B,N],N]=B\), equivalently \([N,C]=-B\), and
\[
 [B,C]e_n=2(b_{n-1}^2-b_n^2)e_n,\qquad b_{-1}=0.
\]
Solving \(k=i(e_--e_+)\) and
\([h,k]=-2i(e_++e_-)\) gives the source's formulas
\(\sigma(e_+)=\tfrac i2(N+c)+\tfrac14C\),
\(\sigma(e_-)=-\tfrac i2(N+c)+\tfrac14C\).
Their commutator is
\(\tfrac i4[N,C]=-iB/4\), proving the defect independent of the
choice of \(b_n\).
For the corrected assignments the commutator is \(i[N,C]=-iB\).
The other two Lie relations are equivalent to
\([B,C]=-4(N+c)\), or
\(b_n^2-b_{n-1}^2=2(n+c)\).  Summing from zero gives precisely
\((n+1)(n+2c)\); substitution verifies sufficiency.
In the actual even oscillator representation the compact generator
has eigenvalues \(2n+1/2\), by S23, equation \texttt{b-k} and
Proposition \texttt{numberN}.  Thus \(2(n+c)=2n+1/2\), so
\(c=1/4\).  This proves the corrected statement without altering
the independently proved archimedean Jacobi coefficients.
\end{proof}

\subsection{Classical prolate spectrum: established facts and remaining identifications}

The self-adjoint full-line prolate operator is not made self-adjoint
by \(A\): that result was already proved by Connes--Moscovici,
\emph{The UV prolate spectrum matches the zeros of zeta},
\emph{PNAS} 119 (2022), e2123174119, Theorem 1.6.  That theorem gives
a discrete spectrum unbounded in both directions, simple negative
eigenvalues, and the extension commuting with both time and Fourier
cutoffs.  Corollary 1.7 places the negative eigenvectors in the Sonin
space.  These are background theorems.
The finite-dimensional possible excess of positive eigenvalues
mentioned in S23 is no longer an open conjecture in the current
literature: Ramis--Richard-Jung--Thomann,
\emph{New eigenfunctions for the negative part of the Connes--Moscovici
prolate spectrum}, \emph{C. R. Math.} 363 (2025), 1065--1081,
Corollary 15, p.~1076, proves that all nonclassical CM eigenvalues are
negative.  This is an unconditional result, not a consequence of \(A\).

The full Darboux square root \(2\!\not\!D\) in CM, Theorem 5.1,
has both real and imaginary spectrum.  It is consequently neither a
self-adjoint nor a skew-adjoint Hilbert-space operator.  A theorem
about the zeros of zeta does not change those existing spectral values.
Its negative-square spectral restriction is meaningful, but identifying
its exact spectral determinant with \(\Xi\) requires an equality
of spectral data, beyond a counting asymptotic.  The same applies to
the semilocal prolate candidates: the required domain, cutoff
commutation, and trace identification have not been supplied by \(A\).

\begin{proposition}[AC-S12: the imaginary-axis differential expression]
\label{AC-S12}
For the classical differential expression
\[
 W_\lambda=-\partial_x((\lambda^2-x^2)\partial_x)
                                  +(2\pi\lambda x)^2,
\]
put \(x=iy\).  The negative of the resulting real-axis expression
is
\[
 B_\lambda=-\partial_y((\lambda^2+y^2)\partial_y)
                                      +(2\pi\lambda y)^2.
\]
Under \(y=\lambda\sinh t\) and the unitary change
\(v(t)=\sqrt{\lambda\cosh t}\,f(\lambda\sinh t)\), it becomes
\[
 -\partial_t^2+(2\pi\lambda^2)^2\sinh^2t
                    +\frac14+\frac{1}{4\cosh^2t}.
\]
On a half-line, parity at zero gives the Dirichlet or Neumann boundary
condition respectively.  In particular its ordinary self-adjoint
realization is strictly positive and has compact resolvent.
\end{proposition}
\begin{proof}
The substitution \(\partial_x=-i\partial_y\) gives the first
formula.  For \(p(y)=\lambda^2+y^2\), the Liouville change
\(dt/dy=p^{-1/2}\), \(v=p^{1/4}f\), gives additional potential
\(p''/4-(p')^2/(16p)\).
Here this is
\(1/2-y^2/(4(\lambda^2+y^2))
 =1/4+1/(4\cosh^2t)\), proving the displayed expression.
The norm is preserved because \(dy=\lambda\cosh t\,dt\).
At zero the multiplier is even with derivative zero, so odd/even
parity gives the claimed separated boundary conditions.
The potential is at least \(1/4\) and tends to infinity at infinity.
The positive closed form is the sum of derivative energy and potential
energy; on a bounded interval its bounded-energy sets are precompact,
while the diverging potential makes the tails uniformly small.
Thus the form embedding is compact, proving compact resolvent and
strict positivity.  This is a differential-expression calculation;
transporting a particular CM boundary condition through analytic
continuation is the separate theorem in RRJT, Theorem 14.
\end{proof}

\subsection{Algebraic-geometric interpretation and its limits}

The arithmetic L-function here is that of the unit Artin--Tate object
over \(\operatorname{Spec}\mathbb Z\).  Its finite-place local factor
comes from the rank-one constant representation and is
\((1-p^{-s})^{-1}\).  The sheaf in AC-S01 lives on the scaling topos,
not on the small etale site of an ordinary scheme.  Its local jet
algebra is a nonreduced zero-dimensional analytic algebra supported
at \(\gamma\).  Passing to the associated graded of its maximal-ideal
filtration gives \(m_\gamma\) copies of its residue field; this is
the semisimplified multiplicity representation used by AC-S08.
Taking this associated graded preserves the local length (and its
class in the Grothendieck group of finite-length modules) but removes
the Jordan extensions exhibited in AC-S02.  No weight or purity
theorem for finite-field etale cohomology removes these extensions
on the basis of \(A\) alone.

AC-S03--AC-S07 are analytic inequalities and convergence statements
for a closed form; they do not construct an algebraic square or an
algebraic correspondence.  AC-S09--AC-S12 concern real/complex
operator domains and an enveloping-algebra action.  Their rigorous
algebraic content is the explicit \(\mathfrak{sl}_2\) module and its
normalization, not a new pure motive or an etale Frobenius action.
Exact archive cotes or scan numbers for these new analytic statements
have not been supplied or verified; none is invented here.
